% Generated draft: requires mathematical and source review.
\documentclass[11pt]{article}
\usepackage[margin=25mm]{geometry}
\usepackage{fontspec}
\setmainfont{Latin Modern Roman}
\usepackage{amsmath,amssymb,mathtools,unicode-math}
\setmathfont{Latin Modern Math}
\usepackage{enumitem,needspace,xurl,hyperref}
\setlength{\parindent}{0pt}
\setlength{\parskip}{5pt}
\setlength{\emergencystretch}{3em}
\newcommand{\N}{\mathbb{N}}
\newcommand{\Z}{\mathbb{Z}}
\newcommand{\Q}{\mathbb{Q}}
\newcommand{\R}{\mathbb{R}}
\newcommand{\C}{\mathbb{C}}
\newcommand{\F}{\mathbb{F}}
\newcommand{\A}{\mathbb{A}}
\newcommand{\PP}{\mathbb{P}}
\newcommand{\E}{\mathbb{E}}
\newcommand{\eps}{\varepsilon}
\newcommand{\dd}{\,\mathrm{d}}
\begin{document}
\section*{UnsolvedMath ID: 20000813}
\textbf{Source problem code:} catalogue entry
\subsection{Definitions and mathematical statement}
In the standard space-curve interpretation, let \(k\) be algebraically closed, \(R=k[x,y,z,t]\) with its standard grading, and let \(C\subset\PP^3_k\) be a locally Cohen--Macaulay curve with ideal sheaf \(\mathcal I_C\). Its Rao module is the finite-length graded module
\[
M_C=\bigoplus_{j\in\Z}H^1(\PP^3,\mathcal I_C(j)).
\]
The source asks what geometric properties of \(C\) follow if \(M_C\) is Gorenstein or annihilated by a nonzero linear form. It remarks that \(C\) might have to be minimal. The ambient curve hypotheses, precise module meaning of Gorenstein, and meaning of minimal are not supplied in the original item. For a cyclic module \(R/J\), Artinian Gorenstein means its socle is one-dimensional; for arbitrary modules one must specify a corresponding duality condition. Minimality is naturally considered in an even liaison class here, but is not a fixed source definition.

The attached aid gives a cohomological consequence. Let \(\ell\in R_1\setminus\{0\}\) annihilate \(M_C\), and require that \(H=V(\ell)\) meets \(C\) properly. Put \(Z=C\cap H\). Define the degree-\(m\) equation-lifting defect by
\[
Q_m(C,H)=
\frac{H^0(H,\mathcal I_{Z,H}(m))}
{\operatorname{im}(H^0(\PP^3,\mathcal I_C(m)))}.
\]
The aid reports the exact identification \(Q_m(C,H)\cong(M_C)_{m-1}\): every Rao class is an obstruction to lifting an equation of the hyperplane section. Proper intersection is essential to the restriction sequence being used; linear annihilation alone does not supply it.

If a chosen graded self-duality \(M_C\cong M_C^\vee(-s)\) is given, it reports perfect pairings
\[
Q_m(C,H)\times Q_{s+2-m}(C,H)\longrightarrow k.
\]
Here \((M^\vee)_j=\operatorname{Hom}_k(M_{-j},k)\) and \(M(a)_j=M_{j+a}\), fixing the shift convention. The pairing depends on the chosen self-duality, not merely an unspecified use of the word Gorenstein.

The aid also gives the explicit cyclic module
\[
M=R/(x,yz,yt,zt,y^4-z^4,y^4-t^4).
\]
It reports that this is an Artinian Gorenstein algebra annihilated by \(x\), but its six minimal first relations violate the applicable four-relation condition for cyclic Rao modules of curves in a double plane. A double plane means the nonreduced subscheme \(V(\ell^2)\), so containment there is much stronger than \(\ell M_C=0\). Invoking Rao's realization theorem, the aid reports an even liaison class whose Rao module is a grading shift of \(M\), with no representative contained in any double plane, including its minimal curves. The full double-plane criterion and citations are not reproduced.

Thus the supplied example rules out this particular proposed geometric conclusion from the two module properties, even after a minimality restriction in that liaison sense. It does not show that the module properties have no useful consequences or exhaust the original request for ``nice properties.'' All these are unreviewed supplied partial claims.
\subsection{Short English statement}
Ask how Gorenstein or linearly annihilated Rao modules constrain space curves, possibly after minimality. The supplied aid identifies hyperplane equation-lifting defects and their duality, while giving a module whose liaison class has no double-plane curve despite both stated algebraic properties.
\subsection{Sources}
\begin{itemize}
\item AIM Problem 12, \emph{Components of Hilbert schemes}, including the remark that the curve may need to be minimal: \url{https://aimath.org/WWN/hilbertschemes/hilbertschemes.pdf}.
\item Local registry \texttt{data/unsolvedmath/problems.json}, permanent ID \(20000813\), exact source code \texttt{catalogue entry}; canonical \texttt{aim-arithmetic-geometry-notes.json}, notes[58].
\item The attached aid invokes Rao's realization theorem and a double-plane relation criterion without complete bibliography, and supplies the defect-pairing and module example. Attempt 1 is explicitly machine-generated, not peer reviewed, and requires independent verification before citation.
\end{itemize}
\end{document}
